No exemplo do misturador de manteiga de amendoim, o texto estabelece que a força de arrasto ($F_D$) depende da velocidade ($V$), mas também é influenciada por três outras variáveis independentes: a largura da lâmina ($d$), a densidade do fluido ($\rho$) e a viscosidade do fluido ($\mu$) [1]. A equação funcional que descreve essa dependência é formulada como:
$$ F_D = F_D(V; d, \rho, \mu) $$

O objetivo experimental hipotético é descobrir como a força de arrasto ($F_D$) varia em função da velocidade ($V$) [1]. Como existem outras três variáveis no sistema, qualquer gráfico bidimensional de $F_D$ versus $V$ precisará manter certos parâmetros fixos. Se escolhermos analisar o sistema usando apenas três (3) valores distintos para cada uma das variáveis extras ($d$, $\rho$ e $\mu$), a combinatória de testes se estrutura da seguinte maneira [2]:

\begin{itemize}
    \item \textbf{Os 9 Gráficos:} Se fixarmos as características do fluido, teremos 3 valores possíveis para a densidade ($\rho$) multiplicados por 3 valores possíveis para a viscosidade ($\mu$). Isso resulta em $3 \times 3 = 9$ misturas ou condições de fluido completamente diferentes [2]. Para cada uma dessas 9 condições de fluido, precisaremos traçar um sistema de eixos independentes relacionando Força versus Velocidade, totalizando assim 9 gráficos diferentes [2].
    
    \item \textbf{As 3 Curvas:} Dentro de cada um desses 9 gráficos (onde $\rho$ e $\mu$ estão fixos para aquele gráfico específico), ainda precisaremos variar a largura da lâmina ($d$). Como escolhemos 3 larguras diferentes, teremos que traçar uma linha (curva) de Força versus Velocidade para a lâmina $d_1$, outra para a lâmina $d_2$ e uma terceira para a lâmina $d_3$ [2]. Isso resulta em exatas 3 curvas dentro de cada gráfico [2].
\end{itemize}

Esta demonstração ilustra o formidável acúmulo de dados experimentais necessários quando se analisa um problema físico sem o auxílio prévio da análise dimensional, que mais tarde provará que toda essa relação complexa pode ser comprimida em grupos adimensionais correlacionados por apenas uma única curva [2].